\section{Introduction}

Recent years have seen rapid progress in large language model agents and increasingly complex multi-agent ecosystems.
Protocols such as the Model Context Protocol (MCP) and Agent-to-Agent (A2A) promise interoperable tool use and agent communication.
In practice, however, current frameworks still tend to expose too much information too early.
Semantic capability descriptions, protocol formats, and execution details are often injected into the agent context upon initialization.

We argue that this is an architectural problem.
Classical agent literature characterizes agents as autonomous entities that act based on their own decision-making process\cite{agent_autonomy_definition}.
Recent work such as MCP-Zero\cite{fei2025mcpzero} similarly argues that exposing full tool schemas to the model creates context overhead and reduces decision autonomy.
In heterogeneous ecosystems, this problem becomes even more pronounced.
Agents must reason not only about which capability they need, but also about protocol-specific execution patterns for both tools and peer agents.
This increases context bloat, raises the risk of protocol confusion, and can degrade end-to-end performance.

To address this, we propose a unified capability gateway that mediates both MCP tools and A2A peer agents through a common capability layer.
Its core purpose is staged information exposure, so that agents can plan on semantic capability information before execution-specific protocol details enter the context.
At the same time, the gateway also acts as a central configuration and governance point, where routing, protocol integration, and cross-cutting control mechanisms can be handled consistently.

We evaluate this idea in a realistic vulnerability-detection multi-agent system.
We compare two direct framework-based baselines against two gateway-based architectures and analyze task success, latency, token consumption, and common failure patterns.
We further study retrieval-based capability selection under growing capability load.
Our results show that gateway-based architectures are more robust and efficient in the evaluated setting, while retrieval becomes especially useful as the surrounding environment grows.

Our main contributions are:
\begin{itemize}
    \item We propose a unified capability gateway for heterogeneous multi-agent systems that enables staged information exposure across planning and execution.
    \item We show empirically that gateway-based architectures outperform direct protocol exposure baselines in robustness, token efficiency, and end-to-end task success in our benchmark setting.
    \item We show that retrieval-based capability selection becomes most useful as the environment grows, as it enables agents to fetch only relevant information.
\end{itemize}
